% ECONOMICS_PROBLEM: {"id": "D91-627", "title": "Present bias and self-control", "jel_code": "D91", "source_id": "OP-0063"}
% !TeX program = xelatex
\documentclass[11pt,a4paper]{article}
\usepackage[margin=23mm]{geometry}
\usepackage{fontspec}
\setmainfont{Latin Modern Roman}
\usepackage{amsmath,amssymb,mathtools}
\usepackage{xcolor,enumitem,booktabs,longtable,array,xurl,hyperref}
\definecolor{muted}{HTML}{53636C}
\hypersetup{colorlinks=true,urlcolor=blue,linkcolor=blue}
\setlength{\parindent}{0pt}
\setlength{\parskip}{5pt}
\newcommand{\R}{\mathbb R}
\newcommand{\E}{\mathbb E}
\newcommand{\N}{\mathbb N}
\newcommand{\ind}{\mathbf 1}
\newcommand{\Var}{\operatorname{Var}}
\newcommand{\Cov}{\operatorname{Cov}}
\newcommand{\supp}{\operatorname{supp}}
\DeclareMathOperator*{\argmin}{arg\,min}
\DeclareMathOperator*{\argmax}{arg\,max}
\newcommand{\entrymeta}[3]{{\sffamily\small\color{muted}\textbf{\texttt{#1}}\quad|\quad #2\par #3\par}}
\newcommand{\Problemref}[1]{\texttt{#1}}
\newcommand{\JELref}[2]{\texttt{#2} (JEL \texttt{#1})}
\begin{document}
\section*{Present bias and self-control}
\textbf{Problem ID:} \texttt{D91-627}\quad \textbf{JEL:} \texttt{D91}\par
% BEGIN ECONOMICS STATEMENT
\label{problem:OP-0063}
\label{atlas:prob:B3}

\noindent\textbf{Original identifier:} B3 (atlas item 63). \textbf{Classification:} Dynamic preference identification research program. \textbf{Original tractability label:} Medium (editorial, not a complexity result).

\subsection*{Definitions and assumptions}

At date $t$, a consumer chooses a consumption and effort plan subject to specified asset, credit, income, and survival constraints. Conditional on information $\mathcal I_t$, use the quasi-hyperbolic objective
\[
V_t=u(c_t,e_t)+\beta\sum_{k=1}^{T-t}\delta^k\mathbb E[u(c_{t+k},e_{t+k})\mid\mathcal I_t],
\]
with $0<\beta\leq1$ and $0<\delta<1$. Future selves reoptimize. Specify whether the consumer correctly anticipates future $\beta$, and specify commitment technologies and their costs. Utility is over consumption and effort, not automatically over dated monetary transfers. Let $\nu$ collect curvature, liquidity, risk, payment reliability, and sophistication parameters.

\subsection*{Formal research question}

For a joint design $D$ combining randomized dated rewards, real-effort tasks, liquidity changes, and commitment offers, characterize the identified set
\[
\mathcal B(P_D)=\{\beta:\exists\delta,\nu\text{ such that }P_D=P_{D,\beta,\delta,\nu}\}.
\]
Determine sufficient restrictions and experimental variation for this set to be a singleton, and construct honest uncertainty regions otherwise. Then test the maintained stability restriction by asking whether the same identified preference parameters predict administrative consumption or effort choices outside the elicitation task. Specify population, follow-up horizon, and prediction-loss tolerance in advance.

\subsection*{Interpretation and scope}

Money arriving earlier may relax a borrowing constraint even when preferences are dynamically consistent. Differential payment risk, transaction costs, or nonlinear consumption responses can likewise mimic a present-bias parameter. Consequently, a comparison of immediate and delayed money identifies $\beta$ only in a model that rules out these alternatives. Real-effort choices reduce some monetary confounds but introduce time-varying effort costs and beliefs about task difficulty. Commitment demand is informative only after excluding ordinary insurance, forced saving, or transaction-cost motives. Repeated elicitation should distinguish parameter drift from measurement error and attrition. For policy evaluation, fix the normative evaluator explicitly: for example, the date-zero self's expected lifetime utility under a chosen commitment opportunity. A different evaluator, such as the contemporaneous self, can rank the same intervention differently; no welfare ordering follows merely from estimating $\beta<1$.

\subsection*{Source audit and status}

[S1] and [S2] provide dynamic models; [S3] studies real-effort inconsistency. The original link [37] resolves to [S4], a lifecycle-estimation working paper originally issued in 2007 whose current record has revised authorship. Its issue year and present author list must not be represented as an unchanged 2007 version. The additional joint identification and validation target is editorial and model-dependent, rather than an unanswered claim that present bias can never be measured.

\subsection*{References}

\begin{enumerate}[label={[S\arabic*]},leftmargin=*]

\item David Laibson (1997). \emph{Golden Eggs and Hyperbolic Discounting}. Quarterly Journal of Economics 112(2), 443--478. \url{https://doi.org/10.1162/003355397555253}\par \textit{Source role:} Dynamic self-control model. \textit{Locator:} Abstract and article metadata.

\item Ted O'Donoghue and Matthew Rabin (1999). \emph{Doing It Now or Later}. American Economic Review 89(1), 103--124. \url{https://www.aeaweb.org/articles?id=10.1257/aer.89.1.103}\par \textit{Source role:} Model of present bias and sophistication. \textit{Locator:} Abstract and article metadata.

\item Ned Augenblick, Muriel Niederle, and Charles Sprenger (2015). \emph{Working over Time: Dynamic Inconsistency in Real Effort Tasks}. Quarterly Journal of Economics 130(3), 1067--1115. \url{https://doi.org/10.1093/qje/qjv020}\par \textit{Source role:} Experimental evidence from effort rather than money. \textit{Locator:} Abstract and article metadata.

\item David Laibson, Sean Chanwook Lee, Peter Maxted, Andrea Repetto, and Jeremy Tobacman (2007). \emph{Estimating Discount Functions with Consumption Choices over the Lifecycle}. NBER Working Paper 13314; current revised author list, original issue year 2007. \url{https://www.nber.org/papers/w13314}\par \textit{Source role:} Linked lifecycle estimation example; version-sensitive authorship. \textit{Locator:} NBER current record and abstract; original issue year versus revised author list.

\end{enumerate}

\subsection*{Common conventions: Source mathematical conventions}

Unless an entry states otherwise, agent and firm sets are finite. State and action spaces are measurable, strategies and outcome maps are measurable, probability laws are specified, and expectations exist for the targets under discussion. Infinite-horizon discount factors lie in $(0,1)$ and discounted objectives are finite. Norms, tolerances and complexity input sizes must be specified for computational comparisons. Constants or primitives that appear locally are inputs, not empirical facts asserted by this catalogue.

\textbf{Identification.} A statistical model $m\in\mathcal M$ implies an observable law $P_m$ and target $\theta(m)$. At law $P$, its identified set is
\[
\Theta(P;\mathcal M)=\{\theta(m):m\in\mathcal M,\ P_m=P\}.
\]
Point identification holds when this set is a singleton, or equivalently when $P_{m_1}=P_{m_2}$ implies $\theta(m_1)=\theta(m_2)$ for all compatible models. Sharp bounds mean bounds attained, or approached when the boundary is not attained, by compatible models. Writing this set is a definition, not a solution: the substantive task is to establish informative restrictions and derive or estimate the set. An unrestricted-confounding model may give only trivial bounds.

\textbf{Inference.} Identification, estimability and valid uncertainty quantification are separate questions. Fix $\alpha\in(0,1)$. A confidence procedure $C_n$ over a declared law class $\mathcal P$ has asymptotic uniform coverage at level $1-\alpha$ when
\[
\liminf_{n\to\infty}\inf_{P\in\mathcal P}\Pr_P\{\theta(P)\in C_n\}\geq1-\alpha.
\]
For set-identified targets the coverage event must be defined separately, for example coverage of the full identified set. If an entry uses identification-indexed models instead of uniquely indexed laws, the same coverage convention is applied over those models. Pointwise limits do not establish uniform validity, and asymptotic validity does not guarantee finite-sample precision. Sample sizes, dependence structures and weak-identification sequences are part of each application.

\textbf{Counterfactuals.} $Y_i(a)$ is a potential outcome under a specified intervention, including its timing, implementation and versions. It is not $Y_i$ conditional on voluntarily choosing $a$. A full intervention vector permits interference; shorthand scalar treatment notation does not silently impose no interference. Assignment, selection, attrition, anticipation and measurement must be specified. A claim about an instrument's local effect is not automatically a population policy effect.

\textbf{Economic equilibrium.} A counterfactual model includes best responses, constraints, feasibility, market clearing, state transitions and expectations consistent with the stated information. When several equilibria exist, either a declared selector is an input or the target is set-valued. Comparisons across policy regimes must use the stated selection convention. Reduced-form relationships are not presumed invariant after an equilibrium-changing intervention.

\textbf{Optimization and welfare.} Optimization is over the stated feasible set. If attainment is not established, $\sup$ or $\inf$ and $\varepsilon$-optimal choices replace an asserted maximizer or minimizer. For example, with value $V=\sup_{a\in\mathcal A}W(a)<\infty$, an $\varepsilon$-optimal set is $\{a\in\mathcal A:W(a)\geq V-\varepsilon\}$ for $\varepsilon>0$. Benefits, costs, risk and health or environmental outcomes can be aggregated only using declared units and conversion weights. Interpersonal utility weights, the treatment of fiscal transfers and foreign profits, discounting, and the behavioral welfare criterion are explicit inputs. A comparative-static derivative additionally requires existence and appropriate differentiability of the selected counterfactual.

\textbf{Sources.} Local labels [S1], [S2], and so forth restart in every question. Locators identify the relevant abstract, model, section or result at the level actually checked; a bibliographic or abstract-level verification is not represented as inspection of an entire proof. Working papers are distinguished from later publications and changing versions. Source summaries are paraphrased. All URL checks and editorial formulations refer to the audit date stated above.
% END ECONOMICS STATEMENT
\end{document}
